\chapter{Notation and Symbols}\label{chap:appA_notation}

This appendix is a quick-reference for the mathematical notation used throughout
the book. It is deliberately terse: each entry gives the printed symbol, the
\LaTeX{} source macro defined in the manual preamble, the meaning of the symbol,
and the typical SI units in which it appears. All math in the manual is built
from the macros listed here, so that a symbol always renders the same way and
always means the same thing. If you are reading the source, you will find these
macros defined once in \code{docs/manual/preamble.tex}; chapter authors are asked
to use them rather than redefine local notation.

A few conventions hold book-wide. Vectors are set in bold (for example the
horizontal velocity $\uvec=(\uu,\vv)$); scalars are set in italic. Horizontal
operators act in the $(\lon,\lat)$ plane on the sphere unless a chapter states
otherwise. The material derivative $\Dt{}$ follows a fluid parcel and expands as
$\Dt{q} = \pd{q}{t} + \uvec\!\cdot\!\grad q + \ww\,\pd{q}{z}$.
Where a quantity has both a continuous (textbook) meaning and a discrete
(code) counterpart, the continuous symbol is used in the equations and the
mapping to the implementation is given in the relevant chapter and in
Appendix~\ref{chap:appC_codemap}; numerical values of constants and tunable
parameters are collected in Appendix~\ref{chap:appB_parameters}.

The governing equations that motivate this notation are introduced in
\ref{chap:ch04_ocean_equations} (ocean) and
\ref{chap:ch09_atmosphere_multilevel} (atmosphere); the standard references for
the symbols and their physical meaning are \citep{vallis2017} and, for the
overturning and box-model variables, \citep{stommel1961,marotzke1991}.

% ------------------------------------------------------------------
\section{Calculus and operators}
% ------------------------------------------------------------------

These are the differential operators used in the conservation laws. The
horizontal gradient, divergence, and curl are understood on the sphere of
radius $\Erad$; their discrete spherical forms (with the metric factors
$\cos\lat$ and the polar treatment) are given in
\ref{chap:ch06_ocean_numerics}.

\begin{table}[htbp]
  \centering
  \caption{Calculus and differential operators.}
  \label{tab:appA_operators}
  \begin{tabular}{@{}llp{0.42\linewidth}l@{}}
    \toprule
    Symbol & Macro & Meaning & Acts on \\
    \midrule
    $\pd{q}{t}$        & \code{\textbackslash pd}   & Partial derivative (local rate of change)            & scalar/field \\
    $\pdd{q}{x}$       & \code{\textbackslash pdd}  & Second partial derivative                            & scalar/field \\
    $\dt{q}$           & \code{\textbackslash dt}   & Ordinary (total) time derivative                     & function of $t$ \\
    $\Dt{q}$           & \code{\textbackslash Dt}   & Material (advective) derivative following a parcel    & field \\
    $\grad q$          & \code{\textbackslash grad} & Gradient operator $\boldsymbol{\nabla}$              & scalar $\to$ vector \\
    $\diver\uvec$      & \code{\textbackslash diver}& Divergence $\grad\!\cdot\!\uvec$                     & vector $\to$ scalar \\
    $\curl\uvec$       & \code{\textbackslash curl} & Curl $\grad\!\times\!\uvec$                          & vector $\to$ vector \\
    $\lap q$           & \code{\textbackslash lap}  & Laplacian $\nabla^{2}q$                              & scalar/field \\
    $\dd x$            & \code{\textbackslash dd}   & Differential (integration / line element)            & --- \\
    \bottomrule
  \end{tabular}
\end{table}

% ------------------------------------------------------------------
\section{Fields and state variables}
% ------------------------------------------------------------------

These are the prognostic and diagnostic fields carried in the model state. The
velocity components $(\uu,\vv,\ww)$ are zonal, meridional, and vertical
respectively; potential temperature $\Temp$ and salinity $\Salt$ set the density
$\dens$ through the equation of state, with $\densz$ a constant reference used in
the Boussinesq approximation (see \ref{chap:ch04_ocean_equations}). The
streamfunction $\psib$ is used both for the horizontal barotropic flow and,
reinterpreted as a meridional overturning streamfunction, for the AMOC
diagnostic in \ref{chap:ch08_ocean_overturning}.

\begin{table}[htbp]
  \centering
  \caption{Fields and state variables.}
  \label{tab:appA_fields}
  \begin{tabular}{@{}llp{0.40\linewidth}l@{}}
    \toprule
    Symbol & Macro & Meaning & Typical units \\
    \midrule
    $\uvec$    & \code{\textbackslash uvec}  & Horizontal velocity vector $(\uu,\vv)$        & m\,s$^{-1}$ \\
    $\uu$      & \code{\textbackslash uu}    & Zonal (eastward) velocity                     & m\,s$^{-1}$ \\
    $\vv$      & \code{\textbackslash vv}    & Meridional (northward) velocity               & m\,s$^{-1}$ \\
    $\ww$      & \code{\textbackslash ww}    & Vertical velocity                             & m\,s$^{-1}$ \\
    $\Temp$    & \code{\textbackslash Temp}  & Potential temperature                         & K ($^{\circ}$C) \\
    $\Salt$    & \code{\textbackslash Salt}  & Salinity                                      & psu (g\,kg$^{-1}$) \\
    $\dens$    & \code{\textbackslash dens}  & Density                                        & kg\,m$^{-3}$ \\
    $\densz$   & \code{\textbackslash densz} & Reference (Boussinesq) density                & kg\,m$^{-3}$ \\
    $\pres$    & \code{\textbackslash pres}  & Pressure                                       & Pa \\
    $\psib$    & \code{\textbackslash psib}  & Streamfunction (barotropic / overturning)      & m$^{3}$\,s$^{-1}$ ($\Sv$) \\
    $\vort$    & \code{\textbackslash vort}  & Relative vorticity $\curl\uvec\!\cdot\!\kvec$  & s$^{-1}$ \\
    \bottomrule
  \end{tabular}
\end{table}

\begin{remark}
The overturning streamfunction is reported in Sverdrups
($1\Sv = 10^{6}\,\mathrm{m^{3}\,s^{-1}}$, macro \code{\textbackslash Sv}); the
AMOC strength at $26.5^{\circ}$N is the central scalar diagnostic of
\ref{chap:ch08_ocean_overturning} and \ref{chap:ch17_amoc_tipping}.
\end{remark}

% ------------------------------------------------------------------
\section{Parameters and constants}
% ------------------------------------------------------------------

These are the geometric, rotational, and gravitational quantities that
parameterise the equations. The Coriolis parameter is $\Cor = 2\Om\sin\lat$ and
its meridional gradient is $\beq = \pd{\Cor}{y}$; on the sphere
$y = \Erad\,\lat$, so $\beq = (2\Om/\Erad)\cos\lat$. The latitude $\lat$ and
longitude $\lon$ are the horizontal coordinates, and $\dep$ denotes the ocean
depth or column thickness. Standard numerical values are tabulated in
Appendix~\ref{chap:appB_parameters}.

\begin{table}[htbp]
  \centering
  \caption{Parameters, constants, and coordinates.}
  \label{tab:appA_parameters}
  \begin{tabular}{@{}llp{0.38\linewidth}l@{}}
    \toprule
    Symbol & Macro & Meaning & Typical units / value \\
    \midrule
    $\Cor$   & \code{\textbackslash Cor}   & Coriolis parameter $2\Om\sin\lat$        & s$^{-1}$ \\
    $\beq$   & \code{\textbackslash beq}   & Coriolis gradient $\pd{\Cor}{y}$         & m$^{-1}$\,s$^{-1}$ \\
    $\Erad$  & \code{\textbackslash Erad}  & Earth radius                              & $6.371\times10^{6}$\,m \\
    $\Om$    & \code{\textbackslash Om}    & Earth rotation rate                        & $7.292\times10^{-5}$\,s$^{-1}$ \\
    $\grav$  & \code{\textbackslash grav}  & Gravitational acceleration                 & $9.81$\,m\,s$^{-2}$ \\
    $\lat$   & \code{\textbackslash lat}   & Latitude                                   & rad ($^{\circ}$) \\
    $\lon$   & \code{\textbackslash lon}   & Longitude                                  & rad ($^{\circ}$) \\
    $\dep$   & \code{\textbackslash dep}   & Ocean depth / column thickness             & m \\
    \bottomrule
  \end{tabular}
\end{table}

% ------------------------------------------------------------------
\section{Software notation}
% ------------------------------------------------------------------

A handful of macros keep references to the implementation consistent. Inline
source identifiers are set in a monospace font with the \code{\textbackslash code}
macro; the model itself is \model, written in \jax{} (the array/autodiff
framework, \citealp{jax2018}), with the multi-level atmospheric dynamical core
provided by \dino{}. The smooth rectifier $\softplus(x)=\log(1+e^{x})$ appears
where a positivity constraint must remain differentiable. Pointers from an
equation to the file that implements it use the code-callout macro, for example:

\codref{chronos\_esm/ocean/diagnostics.py}{AMOC streamfunction $\psib$ and the $26.5^{\circ}$N metric.}

\noindent A complete map from symbols and equations to source files is given in
Appendix~\ref{chap:appC_codemap}, and instructions for running the model in
Appendix~\ref{chap:appD_running}.

\begin{remark}
\model{} is research software. Several components documented in this book are
deliberate simplifications (for example the single-level atmosphere of
\ref{chap:ch10_atmosphere_singlelevel} and the box-model view of the overturning
in \ref{chap:ch17_amoc_tipping}). Where a symbol stands for an idealised or
parameterised quantity rather than a fully resolved one, the chapter that
introduces it says so explicitly; this appendix only fixes the notation, not the
fidelity of the physics it denotes.
\end{remark}
